\documentclass[11pt]{article}

\usepackage[margin=1.1in]{geometry}
\usepackage[T1]{fontenc}
\usepackage[utf8]{inputenc}
\usepackage{lmodern}
\usepackage{microtype}
\usepackage{booktabs}
\usepackage{longtable}
\usepackage{amsmath}
\usepackage[numbers,sort&compress]{natbib}
\usepackage{xcolor}
\usepackage[colorlinks=true,linkcolor=blue!60!black,citecolor=blue!60!black,urlcolor=blue!60!black]{hyperref}
\usepackage{caption}
\captionsetup{font=small,labelfont=bf}

% Tables read from gen/, written by ../brands_paper.py from the committed probe outputs.
% Numbers in the text are written out; gen/numbers.tex lists the headline values for checking after a rerun.
% Each paragraph's source is in a comment below it.
% Plan: PAPER-PLAN-2026-10-04.md §1 (private). Main text uses the one-name questions (Name, Choose, Recommend) and the
% two-turn pick; free forms and the one-turn pick are appendix checks.
% Frame (Tapan, 2026-10-14/15): brand preference for a marketing audience. Leaders keep the Name slot and lose the pick;
% the pick forms at DPO from preferences other models set. Monoculture is motivation and implication, not a result.

\makeatletter\newcommand{\rows}[1]{\@@input #1 }\makeatother

\title{\textbf{How Language Models Learn Taste:\\Brand Preference in LLMs}}
\author{Tapan Parikh\\Cornell Tech\\\texttt{tsp53@cornell.edu}}
\date{October 2026}

\begin{document}
\maketitle

\begin{abstract}
\noindent
Asked to name a running shoe brand, language models say Nike. Asked to recommend one, they say Brooks. This paper asks
where models learn that taste. Across 102 models and 42 brand
categories, the brand a model settles on when recommending differs from the brand it names in 26. The field moves
together to a shared second favourite. Open training checkpoints show where this forms. Pretraining already attaches
the second favourite to advice. Supervised fine-tuning teaches the model to answer and leaves the pick alone.
Preference optimization (DPO) teaches the model which answer is better, and the pick moves in all three pipelines, by 11 to
24 points. The preference data contains almost no brands. DPO teaches a general preference for the considered answer over the
obvious one, and it moves other judgments the same way. Asked to choose a sauce, models switch from ketchup to pesto.
Other models wrote or judged the preferred answers.
\end{abstract}
% 102, 42, 26: probes/brand_ladder.json levels_specific (gen/ladder_rows.tex). Labs: brand_confounds.json by_lab.
% Stages: brand_logprob.json steps sft->dpo pick_h2h (gen/stage_pick_rows.tex). Data: stage_dolci.json, stage_tulu.json.
% Sauce: answer_logprob_steps.json. Census: parikh2026census.

\section{Introduction}
\label{sec:intro}

Ask a language model to name a running shoe brand and it says Nike. Ask it to recommend one and it says Brooks. The
first question asks what comes to mind. The second asks for a judgment.
% brand_ladder.json levels_specific running shoe.

Across 102 models, the brand they settle on when recommending differs from the brand they name in 26 of 42
categories. The models converge on the same second favourite. Delta gives way to Singapore Airlines, Chrome to Firefox
and Google to DuckDuckGo. Every major lab makes the move. The named brands are market leaders, and the picks are mostly challengers and
specialists. People increasingly ask assistants what to buy. A leader keeps the slot for what comes to mind and loses
the slot for what to buy, and every assistant steers buyers toward the same alternatives
\citep{kleinberg2021monoculture,bommasani2022picking}.
% gen/ladder_rows.tex, gen/flip_rows.tex; brand_confounds.json by_lab.

This paper asks where models learn this taste. We use open models that release a checkpoint after each stage of
training. Pretraining supplies the candidates. A base model already leans to the second favourite when a sentence is
phrased as advice. Supervised fine-tuning (SFT) teaches the model to answer and leaves the pick alone. Direct
preference optimization (DPO) teaches the model which of two answers is better, and the pick moves at DPO in all three
pipelines. Reinforcement learning (RL) leaves it there.
% gen/stage_pick_rows.tex; advice_register.json.

DPO's data holds almost no brand comparisons. The same step moves judgments that have nothing to do with brands.
Asked to name a sauce, models say ketchup. Asked to choose one, OLMo 3 7B says pesto 5\% of the time as a base model
and 67\% of the time after DPO. For brands, DPO moves recommendations away from the brand the base model finds most
obvious. Other models wrote or judged the preferred answers. A brand's own content cannot reach this preference.
% stage_dolci.json, stage_tulu.json; answer_logprob_steps.json (sauce).

\section{Related work}
\label{sec:related}

\paragraph{Brand recommendations.} Audits of production assistants measure which brands surface and how stable they
are. Paraphrases of one buying intent share few recommendations \citep{jack2026paraphrase}, retrieval moves small brands
more than priors do \citep{tannenbaum2026stage}, and well-known brands win ties on identical specifications
\citep{chu2026brands}. Audits average over wordings \citep{zatuchin2026dice}. Changing the frame of a question changes
model answers \citep{trhlik2026multitudes,rabbani2026fact,tonneau2026cues}. We study one such change across a field of
models and trace it through training.

\paragraph{Convergence and training stages.} Models converge on open-ended outputs \citep{jiang2025hivemind} and on
one-word defaults \citep{parikh2026census}. On the same open pipelines we use, output homogeneity is traced to SFT
\citep{fortier2026convergence} and diversity loss to DPO \citep{karouzos2026where}. On Tulu, preference data matters more
than the optimization algorithm \citep{ivison2024unpacking}.

\paragraph{Transmission between models.} Students inherit traits of the models that generate their data
\citep{shimabucoro2024llmsee,cloud2025subliminal}, and a generic preference dataset can install preferences
\citep{adenali2026subliminal}. Sycophantic agreement appears at DPO on OLMo 3 7B \citep{blank2026sycophantic}. LLM
judges favour students of their own lineage \citep{li2025leakage}. Human raters reward assertive answers
\citep{hosking2023human}.


\section{Setup}
\label{sec:method}

\paragraph{Questions.} The brand battery has 42 categories, from soda to cloud providers. Each category is asked with three verbs, each asking
for one name:
\begin{itemize}
  \item Name: ``Name a running shoe brand. Reply with the name only.''
  \item Choose: ``Choose a running shoe brand. Reply with the name only.''
  \item Recommend: ``Can you recommend a running shoe brand? Reply with the name only.''
\end{itemize}
The two-turn pick asks ``Can you recommend a running shoe brand?'' and then ``Which one would you pick?'' in the same
conversation. Appendix~\ref{app:panel} checks the three verbs without the one-name instruction.
% Wordings: spec/stimulus_brands*.json, spec/perturb/stimulus_brands_*_clamp.json, spec/stimulus_brands_pick.json.

\paragraph{Panel.} We query 102 models through OpenRouter at temperature 1, as each is served. The one-name questions run
8 times per model and category, the two-turn pick 4 times. Twenty-five hybrid models were also run with reasoning turned
off. Runs were collected between 2026-09-30 and 2026-10-10. Perplexity's Sonar answers from live search and is left
out.
% spec/runs.json (brands*, hybrids); brand_ladder.EXCLUDE. Dates from transcript run_date fields.

\paragraph{Scoring.} A one-name reply counts if it is one name, with no list and no lead-in. A two-turn reply is
reduced to the brand it commits to. A list, a question back or a refusal counts as no pick, which happens in 17\% of
two-turn replies. On three hand-checked samples the committed-brand rules agreed with a human reading on 87
to 97\% of replies. Aliases map variant spellings and product names to one brand.
% brand_ladder.py docstring (hand checks 2026-10-04, 120/60/30 replies).

\paragraph{Measures.} For each level we report the number of categories whose consensus brand differs from Name's, the
mean share of the top brand, and how often a model keeps its own one-word Name default among replies that name a
brand. Appendix~\ref{app:panel} reads each contrast against a split-half floor.
% factors.py.

\section{The field names one brand and picks another}
\label{sec:panel}

Table~\ref{tab:ladder} gives each level. Under Name every category has a clear consensus, and
models keep their own default 89\% of the time. Each step toward a judgment moves the consensus in more categories.
Choose moves it in 6 of 42, Recommend in 16, and recommending then picking in 26. Over the same steps a model's own
default falls from 89\% of its answers to 31\% of its picks.

\begin{table}[t]
\centering\small
\caption{From naming to picking, on the panel (102 models, 42 categories). ``Differs'' counts the categories whose
consensus brand differs from Name's. ``Top share'' is the mean share of the consensus brand. ``Kept'' is how often a
model's answer is its own one-word Name default, among replies that name a brand.}
\label{tab:ladder}
\begin{tabular}{lrrr}
\toprule
Question & Differs & Top share & Kept \\
\midrule
\rows{gen/ladder_rows.tex}
\bottomrule
\end{tabular}
\end{table}
% probes/brand_ladder.json levels_specific.

The consensus holds as models leave the default. The top brand's share is 76\% under Name and between 60
and 69\% at every other level. The models move together to a new brand. Table~\ref{tab:flips} lists the 26 categories
where the two-turn pick leaves the Name brand. Nike becomes Brooks, Chrome becomes Firefox, Delta becomes Singapore
Airlines and WhatsApp becomes Signal. Toyota, Coca-Cola, Amazon and LEGO hold at every level.

The named brands are market leaders: Google, Walmart, McDonald's, Facebook, Verizon. The picks are mostly smaller
brands with a specialist or insider following: DuckDuckGo, Aldi, Sierra Nevada, Reddit, Signal, Brooks. A few picks
are premium brands, such as Hermès over Chanel and Hyatt over Marriott.
% gen/ladder_rows.tex (top share column); gen/flip_rows.tex.

\begin{table}[t]
\centering\small
\caption{Consensus brand by level, for the 26 categories where the two-turn pick differs from Name.}
\label{tab:flips}
\setlength{\tabcolsep}{4pt}
\begin{tabular}{lllll}
\toprule
Category & Name & Choose & Recommend & Recommend, then pick \\
\midrule
\rows{gen/flip_rows.tex}
\bottomrule
\end{tabular}
\end{table}
% probes/brand_ladder.json levels_specific consensus.

In the two-turn pick, a model's own default appears in its recommendation list 83\% of the time and is its pick 31\%
of the time. Models list their default and choose something else.
% brand_ladder.json two_turn.

In each of the 12 labs with at least three models, the two-turn pick moves the lab's consensus brand in 55 to 74\% of
categories. Labs favour their own assistant. Anthropic's models recommend Claude 92\% of the time, against 15\% for
other labs' models, which inflates the AI assistant row. Labs do not favour their other products: Google's models
recommend Google search and Chrome less often than other models do. Two halves of the Name runs disagree on the consensus brand in no category. Appendix~\ref{app:panel} gives
this and the other checks.
% brand_confounds.json by_lab (consensus_moved_pick2), paraphrase; probes/factors.json.

These models' training is closed. To find where the taste forms, we turn to open models.

\section{Taste forms when a model learns which answer is better}
\label{sec:stages}

Three pipelines release a checkpoint after each stage of training: OLMo 3 7B, OLMo 3.1 32B and Tulu 3 8B. Each has a
base model, an SFT stage, a DPO stage and an RL stage. Nemotron 3.5 Lightning releases a base and a final model. We read the base stage as text completion and the tuned stages through their chat format with no system prompt. We
score by log-probability, which measures every stage on every category without sampling.
% harness/ladders.json; probe_brand_logprob.py.

\paragraph{Pretraining supplies the second favourite.} A base model cannot be asked for advice, but it can complete an
advice sentence. Given ``If you want a good web browser, I recommend'', Tulu's base model continues with Firefox in 13
of 20 samples. Given ``A well-known web browser is'', it continues with Chrome in 10. Over the 16 categories where the
panel's Recommend brand differs from its Name brand, advice phrasing raises that brand's share in 7 to 10 categories,
depending on the base model, and lowers it in at most one. Log-probabilities agree in every base model (Appendix~\ref{app:stages}).
% probes/advice_register.json; probe_brand_logprob.py stage base (gen/stage_shift_rows.tex).

\paragraph{Learning to answer leaves the pick.} To compare stages on the same choice, we give every stage the same
three-item list and ask ``Which one would you pick?''. The list holds the Name brand, the pick brand and a third brand
the panel recommends. The reply is started with ``I'd pick'', and we read the probability of each brand. Each list is
shown in both orders, because base models favour the brand listed first. Table~\ref{tab:stages} gives the share for
the pick brand over the Name brand. After SFT it sits near the base value in all three pipelines.
% probes/brand_logprob.json stages *.pick_h2h.

\begin{table}[t]
\centering\small
\caption{Fixed-list pick by stage: probability of picking the panel's pick brand over its Name brand, averaged over
both list orders and over the 26 categories where the two differ. Last column: the SFT-to-DPO change with a 90\%
bootstrap interval over categories.}
\label{tab:stages}
\begin{tabular}{lrrrrl}
\toprule
Pipeline & Base & SFT & DPO & RL & SFT $\to$ DPO \\
\midrule
\rows{gen/stage_pick_rows.tex}
\bottomrule
\end{tabular}
\end{table}
% probes/brand_logprob.json stages *.pick_h2h, steps sft->dpo.

\paragraph{Learning the better answer moves it.} DPO raises the pick brand's share in all three pipelines, by 24
points in OLMo 3 7B, 11 in Tulu 3 8B and 16 in OLMo 3.1 32B. RL leaves it where DPO put it. Nemotron, which has only a
base and a final model, rises from 58\% to 79\%.

The move does not depend on which brands the panel picks. Over all 42 categories, we compare each brand's
log-probability under Name at the base stage with how much DPO raises it under Recommend. The correlation is negative
in all three pipelines ($-0.16$, $-0.14$ and $-0.23$, each interval below zero). DPO raises the brands the base model
found less obvious.
% gen/stage_pick_rows.tex; gen/numbers.tex dpoStep*. Nemotron: brand_logprob.json stages nemotron.

\paragraph{The same step on the models' own lists.} We repeat the measure on the SFT stage's own replies to ``Can you recommend\ldots'', keeping those that name both brands (89 to 146 lists per
pipeline). The DPO step remains, at $+14$, $+9$ and $+14$ points. Appendix~\ref{app:stages} compares log-probabilities with sampled picks and checks a known side effect of DPO.
% brand_logprob_ownlist.json.

\section{What preference optimization teaches}
\label{sec:general}

\paragraph{The data holds almost no brands.} OLMo 3's SFT and DPO data and Tulu 3's SFT and preference mixtures are
public. We searched them for requests that ask for a recommendation in the 26 categories and read the first brand the
answer names (Table~\ref{tab:data}). In 260,000 OLMo 3 DPO pairs, the chosen answer names the pick brand 3 times and the rejected answer names it once.
Tulu 3 is as sparse. DPO did not learn these brands from its data.
% probes/stage_dolci.json (dpo_pairs 259,922), stage_tulu.json.

\begin{table}[t]
\centering\small
\caption{Brand recommendations in the post-training data, over the 26 categories: first brand named in SFT answers, and
in the chosen and rejected answers of preference pairs.}
\label{tab:data}
\begin{tabular}{lrrrrrr}
\toprule
 & \multicolumn{2}{c}{SFT} & \multicolumn{2}{c}{DPO chosen} & \multicolumn{2}{c}{DPO rejected} \\
Data & Name & Pick & Name & Pick & Name & Pick \\
\midrule
\rows{gen/data_rows.tex}
\bottomrule
\end{tabular}
\end{table}
% probes/stage_dolci.json, stage_tulu.json.

\paragraph{The same step moves judgments outside brands.} The census asks the same everyday categories with Name and
Choose \citep{parikh2026census}. In 23 of them the panel's Choose answer differs from its Name answer. Asked for a
sauce, the panel names ketchup and chooses pesto. We read the same log-probability shift on these 23 categories
(Table~\ref{tab:census}). Base models barely separate the two verbs. DPO moves Choose toward the panel's answer in all three pipelines.
Outside brands, the step does not favour less obvious answers in general: the same correlation is near zero in all
three pipelines. It moves Choose toward the specific answers the field gives. For OLMo 3 7B asked to choose a sauce, pesto goes from 5\% at base to 27\% after SFT and 67\% after
DPO, and ketchup falls from 72\% to 6\%.
% probes/answer_logprob_steps.json.

\begin{table}[t]
\centering\small
\caption{Census Choose: how far ``Choose'' moves the odds toward the panel's Choose answer, relative to ``Name'' (nats),
over the 23 categories where the two differ. Last column: the SFT-to-DPO change, 90\% interval.}
\label{tab:census}
\begin{tabular}{lrrrrl}
\toprule
Pipeline & Base & SFT & DPO & RL & SFT $\to$ DPO \\
\midrule
\rows{gen/census_rows.tex}
\bottomrule
\end{tabular}
\end{table}
% probes/answer_logprob_steps.json.

\paragraph{Other models define the better answer.} Half of OLMo 3's DPO pairs are judged by an LLM. The other half pair
a response from Qwen 3 32B, as chosen, with one from a much smaller model, as rejected \citep{geng2025delta}. Tulu 3's
pairs are rated by GPT-4o \citep{lambert2024tulu}. In both pipelines other models decide which answer is better, and the student applies what it learns to brands it
never saw compared.
% Dolci preference_type counts (llm_judged 124,980; delta_learning 124,942). Tulu judge: GPT-4o-2024-08-06, rating 1-5 on four aspects (Tulu 3 report, preference data section).

\paragraph{Teacher or own favourite.} Tulu 3's preference data has one judge. Answers from a pool of 22 models are
rated by GPT-4o \citep{lambert2024tulu}, and the top-rated answer becomes the chosen one. Where GPT-4o recommends a different brand from Tulu's
own base-model favourite, we give each stage a list with both, in both orders. DPO moves the pick toward GPT-4o's
brand, by 8 points (interval $-14$ to $-2$). OLMo 3's data has no single teacher, so the test applies to Tulu only (Appendix~\ref{app:stages}).
% probe_brand_teacher.py; probes/brand_teacher.json (teachers: probes/teacher_brands/ for Qwen 3 32B, the panel's
% gpt-4o file for Tulu; the panel's gpt-4o is the current alias, not the 2024-08-06 judge snapshot). 24, 25 and 23 categories where teacher and student brands differ.

\section{Conclusion}
\label{sec:discussion}

Language models name the market leader and recommend a challenger, and the field agrees on which challenger. On
open pipelines the challenger is present after pretraining, attached to advice. SFT leaves the pick alone. DPO makes
the model choose it, using preference data that holds almost no brands and that other models wrote or judged.

For brands this moves the source of a recommendation. Content and retrieval act on what a model knows about a brand.
The preference that decides the pick is set in post-training, attaches to a kind of brand, and is shared across
vendors. A leader keeps the slot for what comes to mind and loses the slot for what to buy. Being seen as the
specialist's choice counts for more than being the best known.

\section{Limitations}
\label{sec:limits}

Our questions are short and state no needs. Real requests state needs, and we do not measure how much the taste
shapes the answer then. Measures on short prompts can fail to predict realistic use \citep{lum2025trick}. Audits of
production assistants measure real requests and are the complement to this work. Assistants that search the web
could also answer differently. In a spot check, Claude Opus 5.5 answered the two-turn question in all 42 categories
with web search available and searched in none of them. It left the leader in 25 of the 26 categories where the panel
does, and kept the panel's brand in 14 of the 16 where the panel holds.
% probes/search_spotcheck.json (one sample per category, run as Claude Code subagents, 2026-10-15).
% Framing: Tapan, 2026-10-14.

The three staged pipelines come from one lab and share much of their recipe. Nemotron agrees from base to final but has
no intermediate stages. The 32B and Nemotron run at 8-bit precision. For OLMo 3 7B, sampled picks place the change at
RL (Appendix~\ref{app:stages}). The panel is a dated specimen of models as served between 2026-09-30 and 2026-10-10.


\bibliographystyle{plainnat}
\bibliography{references}

\clearpage
\appendix

\section{Panel checks}
\label{app:panel}

\paragraph{Floors and the one-name instruction.} Table~\ref{tab:factors} reads each verb contrast against a split-half floor. Two
halves of the Name runs disagree on the consensus brand in no category, while Name and Recommend disagree in 38\% of
categories. Dropping the one-name instruction leaves Name nearly unchanged and moves Choose and Recommend in about a third of
categories. The verb effect holds in both forms.

\begin{table}[h]
\centering\small
\caption{Verb and one-name effects against split-half floors. ``Moved'' is the share of categories whose consensus brand
changes. ``Kept'' is how often a model's modal answer survives the change.}
\label{tab:factors}
\begin{tabular}{lrrrr}
\toprule
Contrast & Moved & Floor & Kept & Floor \\
\midrule
\rows{gen/factor_rows.tex}
\bottomrule
\end{tabular}
\end{table}
% probes/factors.json, battery brands.

\paragraph{Paraphrase.} Rewording Recommend twice, on a 14-model subset, leaves its consensus brand unchanged in 34 to
36 of 39 categories.
% brand_confounds.json paraphrase.

\paragraph{One-turn pick.} The one-turn question ``Which running shoe brand would you pick?'' moves the consensus in
15 of 42 categories with the one-name instruction. Without it, it names no brand in 58\% of replies, so it measures refusal more than
preference.
% brand_ladder.json levels_specific pick1_clamp, pick1.

\section{Stage checks}
\label{app:stages}

\paragraph{The direct answer.} Table~\ref{tab:shift} reads how far ``Can you recommend\ldots'' moves the odds from the
Name brand toward the pick brand, relative to ``Name\ldots''. The shift is already present in every base model. It
rises at DPO in all three pipelines, and two rewordings of Recommend rise with it. OLMo 3.1 32B also shows a large
rise at SFT in the direct answer, on brands and on census categories alike. Its pick from a list does not rise at SFT.

\begin{table}[h]
\centering\small
\caption{How far ``Can you recommend\ldots'' moves the odds from the Name brand toward the pick brand, relative to
``Name\ldots'' (nats), by stage. Last column: the SFT-to-DPO change, 90\% interval.}
\label{tab:shift}
\begin{tabular}{lrrrrl}
\toprule
Pipeline & Base & SFT & DPO & RL & SFT $\to$ DPO \\
\midrule
\rows{gen/stage_shift_rows.tex}
\bottomrule
\end{tabular}
\end{table}
% probes/brand_logprob.json stages *.pick.shift.recommend; answer_logprob_steps.json (census SFT step, 32B).

\paragraph{Sampled picks.} Sampling picks from the SFT stage's own lists agrees with the log-probabilities within a
few points everywhere except OLMo 3 7B at DPO (Table~\ref{tab:agree}). There, sampled picks do not move until RL.
First-token probabilities and generated answers are known to diverge in tuned models \citep{wang2024answerc}.

\begin{table}[h]
\centering\small
\caption{The pick on the SFT stage's own lists: sampled picks (5 per list) against log-probabilities.}
\label{tab:agree}
\begin{tabular}{llrr}
\toprule
Pipeline & Stage & Sampled & Log-probability \\
\midrule
\rows{gen/agreement_rows.tex}
\bottomrule
\end{tabular}
\end{table}
% stage_pick.json held_fixed_pooled (complete); brand_logprob_ownlist.json.

\paragraph{The base runner-up.} DPO can concentrate probability on outputs that were already likely
\citep{ren2024dynamics} or move it to unrelated ones \citep{razin2024likelihood}. The first would move a model toward
its base runner-up. The base model's second choice under Name is the panel's pick brand in only 2 or 3 of 26
categories, and DPO does not favour the runner-up over the pick.
% probes/obvious.json runner_up (probe_obvious.py).

\paragraph{The teacher test on OLMo.} Half of OLMo 3's preference pairs set an answer from Qwen 3 32B against one from
a much smaller model, with no judge, and the other half are judged across a dozen models. Taking Qwen 3 32B's brand as
the teacher's, the OLMo 3 7B and OLMo 3.1 32B pipelines move $-4$ and $+3$ points at DPO, and both intervals span zero.
% probes/brand_teacher.json.

% Deferred from the plan, numbers to pull before including: identity prompts (OLMo self-naming 0/20 -> 19/20, 5 of 37
% brand defaults; cite lehr2025selfpreference) and the 21-language pilot (brands_lang.py).

\end{document}
